\documentclass[11pt]{article}

\usepackage[margin=1in]{geometry}
\usepackage{amsmath,amssymb,amsthm,mathtools}
\usepackage{graphicx}
\usepackage[hidelinks]{hyperref}

\newtheorem{theorem}{Theorem}
\newtheorem{lemma}[theorem]{Lemma}
\newtheorem{proposition}[theorem]{Proposition}
\newtheorem{corollary}[theorem]{Corollary}

\newcommand{\CR}{\operatorname{CR}}
\newcommand{\DR}{\operatorname{DR}}
\newcommand{\N}{\mathbb{N}}
\newcommand{\eps}{\varepsilon}
\newcommand{\directlim}{\varinjlim}
\newcommand{\Sphere}{S}

\title{A Counterexample to the Finite-Complexity-Rate Inductive-Limit Question}
\author{arXiv automatic math run \texttt{fa\_banach\_001}}
\date{June 22, 2026}

\begin{document}
\maketitle

\section*{Status}

This packet gives a candidate counterexample, likely valid, to the inductive-limit
question in Tian--Yu \cite{TianYu2026}.  The result is stronger than the earlier
partial packet: it gives a negative answer both for injective/direct-union
inductive limits and for filtered colimits with arbitrary bonding
homomorphisms.

The supplied report
\[
\texttt{finite\_complexity\_rate\_counterexample.tex/pdf}
\]
has been copied into the evidence folder.  The previous direct-union reduction
packet has been absorbed into the history folder of this packet.

\section*{Source Question}

The source is Geng Tian and Guoliang Yu, \emph{Embedding complexity into Banach
spaces and the strong Novikov conjecture}, arXiv:2605.12930v1.  After
Definition 1.4 and Proposition 1.5, the paper asks on page 6:

\begin{quote}
``It is natural to ask whether every countable discrete group is an inductive
limit of groups with finite complexity rates.''
\end{quote}

\begin{center}
\includegraphics[width=\textwidth]{figures/open_problem_crop.png}
\end{center}

\section*{Definitions Used}

Let \(\Gamma\) be a countable group with a left-invariant proper metric.  In
Definition 1.4 of \cite{TianYu2026}, finite complexity rate means
\(\CR(\Gamma)<\boldsymbol{\infty}\), equivalently \(\DR(\Gamma)>1\).  Thus
there is an element
\[
  \mathbf a=(a_0,a_1,\ldots,a_k)>1
\]
in the lexicographic order used in the source, a constant \(C>0\), thresholds
\(N_p>0\), and maps
\[
  \phi_p:\Gamma\longrightarrow \Sphere(\ell^{2p}(\N))
\]
such that
\begin{align}
 d(x,y)\le 1
 &\Longrightarrow
 \|\phi_p(x)-\phi_p(y)\|_{2p}\le \eps_p, \label{eq:variation}\\
 d(x,y)\ge N_p
 &\Longrightarrow
 \|\phi_p(x)-\phi_p(y)\|_{2p}\ge 1, \label{eq:separation}
\end{align}
where
\[
 \eps_p=
 \frac{C}{p^{a_0}(\log p)^{a_1}\cdots(\log_k p)^{a_k}}.
\]
The source's truncated convention for small iterated logarithms is irrelevant
for the asymptotic argument.  Since \(\mathbf a>1\), one has
\[
  p\eps_p\longrightarrow 0. \tag{1}
\]

\section*{Proof Intuition}

Finite complexity rate would provide almost unit-scale-constant maps
\(\phi_p\) into the unit spheres of \(\ell^{2p}\), with variations
\(\eps_p=o(1/p)\), while still separating sufficiently far group elements.
Expanders forbid this rate.  If an expander family coarsely embeds in a
finitely generated group, then sufficiently many pairs in a large graph are
sent beyond the separation threshold.  The \(L_q\) Poincare inequality for
expanders, with \(q=2p\), forces the edge variation of the composed map to be
at least a constant multiple of \(1/p\).  This contradicts \(p\eps_p\to 0\).
Osajda supplies a finitely presented group with exactly this coarse expander
containment, giving the counterexample.

\section*{Direct-Union Reduction}

\begin{lemma}[Subgroup heredity]\label{lem:subgroup}
If a countable group \(\Gamma\) has finite complexity rate, then every subgroup
\(\Lambda\leq\Gamma\) has finite complexity rate.
\end{lemma}

\begin{proof}
Restrict a left-invariant proper metric on \(\Gamma\) to \(\Lambda\).  It is
again left-invariant and proper.  Restricting the maps \(\phi_p\) preserves the
variation and separation estimates, and \cite{TianYu2026} states that the
definition is independent of the chosen left-invariant proper metric.
\end{proof}

\begin{proposition}[Directed-union reduction]\label{prop:directunion}
For a countable group \(\Gamma\), the following are equivalent:
\[
\Gamma \text{ is a directed union of finite-complexity-rate subgroups}
\]
and
\[
\text{every finitely generated subgroup of } \Gamma
\text{ has finite complexity rate.}
\]
\end{proposition}

\begin{proof}
If \(\Gamma=\bigcup_{i\in I}\Gamma_i\), with \(I\) directed and all
\(\Gamma_i\) of finite complexity rate, then any finitely generated subgroup
\(H\leq\Gamma\) has a finite generating set contained in one stage
\(\Gamma_i\).  Lemma~\ref{lem:subgroup} gives finite complexity rate for
\(H\).

Conversely, the finitely generated subgroups of \(\Gamma\), ordered by
inclusion, form a directed system whose union is \(\Gamma\).  If each of them
has finite complexity rate, this is the required directed union.
\end{proof}

\section*{The Expander Obstruction}

Let \(X=(V,E)\) be a finite connected \(d\)-regular graph, and let \(A_X\) be
the normalized adjacency operator.  A sequence \((X_n)\) is a spectral expander
sequence with gap \(\sigma>0\) if \(|V(X_n)|\to\infty\) and
\[
  1-\lambda_2(A_{X_n})\ge \sigma
\]
for every \(n\).

We say that \((X_n)\) coarsely embeds into a finitely generated group
\(\Gamma\) if there are maps \(\iota_n:V(X_n)\to\Gamma\) and nondecreasing
functions \(\rho_-,\rho_+:[0,\infty)\to[0,\infty)\), independent of \(n\), with
\(\rho_-(t)\to\infty\), such that
\[
 \rho_-\bigl(d_{X_n}(u,v)\bigr)
 \le d_\Gamma\bigl(\iota_n(u),\iota_n(v)\bigr)
 \le \rho_+\bigl(d_{X_n}(u,v)\bigr).
\]

\begin{lemma}[\(L_q\) Poincare inequality]\label{lem:poincare}
Let \(X=(V,E)\) be a finite \(d\)-regular graph whose normalized adjacency
operator has spectral gap at least \(\sigma>0\).  For every \(q\ge 2\) and
every map \(f:V\to\ell^q\),
\[
 \left(
   \frac{1}{|V|^2}\sum_{u,v\in V}\|f(u)-f(v)\|_q^q
 \right)^{1/q}
 \le
 \frac{2q}{\sqrt\sigma}
 \left(
   \frac{1}{d|V|}\sum_{u\in V}\sum_{v\sim u}
        \|f(u)-f(v)\|_q^q
 \right)^{1/q}.
\]
\end{lemma}

\begin{proof}
This is the \(L_q\) expander Poincare inequality of Matousek, in the form
quoted as Corollary 4.5 in Naor--Silberman \cite{NaorSilberman2011}.
\end{proof}

\begin{theorem}[Expanders force the critical \(1/p\) rate]\label{thm:obstruction}
Let \(\Gamma\) be finitely generated.  Suppose a bounded-degree spectral
expander sequence with gap \(\sigma>0\) coarsely embeds into \(\Gamma\).  If
maps \(\phi_p:\Gamma\to \Sphere(\ell^{2p}(\N))\) satisfy
\eqref{eq:variation} and \eqref{eq:separation}, then there is a constant
\(c>0\), independent of \(p\), such that
\[
  p\eps_p\ge c \qquad (p\in\N).
\]
Consequently \(\Gamma\) has infinite complexity rate.
\end{theorem}

\begin{proof}
Use a word metric on \(\Gamma\), and set
\[
  L=\max\{1,\lceil \rho_+(1)\rceil\}.
\]
Since the word metric is geodesic, \eqref{eq:variation} implies
\[
 d_\Gamma(x,y)\le L
 \Longrightarrow
 \|\phi_p(x)-\phi_p(y)\|_{2p}\le L\eps_p. \tag{2}
\]

Fix \(p\), set \(q=2p\), and choose \(R_p\) with
\(\rho_-(R_p)\ge N_p\).  Since the degrees are bounded, every radius
\((R_p-1)\)-ball in the graphs has size at most a number \(B_p\) independent
of \(n\).  Choose \(n\) with \(|V(X_n)|\ge 2B_p\).  Then at least half of the
ordered pairs \((u,v)\in V(X_n)^2\) satisfy \(d_{X_n}(u,v)\ge R_p\).

Let \(f=\phi_p\circ \iota_n\).  For those far pairs, the coarse lower control
and \eqref{eq:separation} give \(\|f(u)-f(v)\|_q\ge 1\).  Hence
\[
 \left(
   \frac{1}{|V(X_n)|^2}
   \sum_{u,v\in V(X_n)}\|f(u)-f(v)\|_q^q
 \right)^{1/q}
 \ge 2^{-1/q}. \tag{3}
\]
For adjacent vertices \(u\sim v\), coarse upper control gives
\(d_\Gamma(\iota_n(u),\iota_n(v))\le \rho_+(1)\le L\), so by (2)
\[
  \|f(u)-f(v)\|_q\le L\eps_p.
\]
Applying Lemma~\ref{lem:poincare} and (3),
\[
  2^{-1/q}\le \frac{2q}{\sqrt\sigma}L\eps_p.
\]
With \(q=2p\), this gives
\[
  p\eps_p\ge
  \frac{\sqrt\sigma}{4L}2^{-1/(2p)}
  \ge \frac{\sqrt\sigma}{4\sqrt2\,L}.
\]
This lower bound contradicts \(p\eps_p\to 0\) whenever \(\Gamma\) has finite
complexity rate.
\end{proof}

\section*{The Counterexample}

Osajda constructed finitely generated groups whose Cayley graphs contain
isometrically embedded expander sequences \cite[Theorem 3.2 and Corollary
3.3]{Osajda2020}.  More strongly for the present purpose, he obtained closed
aspherical manifolds of dimension at least four whose fundamental groups
contain coarsely embedded expanders \cite[Corollary 3.5]{Osajda2020}.  Such a
fundamental group is finitely presented.

\begin{corollary}\label{cor:fpbad}
There exists a countable finitely presented group \(G\) with
\(\CR(G)=\infty\).
\end{corollary}

\begin{proof}
Take \(G\) to be one of Osajda's finitely presented fundamental groups
coarsely containing a bounded-degree expander family.  Cheeger's inequality
turns the uniform combinatorial expansion in the construction into a uniform
spectral gap after passing to a fixed-degree regular expander family.  Theorem
\ref{thm:obstruction} gives \(\CR(G)=\infty\).
\end{proof}

\begin{lemma}[Finitely presented colimits retract from a stage]\label{lem:retract}
Let \(G\cong\directlim_i G_i\) be a filtered colimit of groups.  If \(G\) is
finitely presented, then for some \(i\) there is a homomorphism
\(s:G\to G_i\) such that the canonical map \(\lambda_i:G_i\to G\) satisfies
\(\lambda_i\circ s=\operatorname{id}_G\).  In particular, \(G\) is isomorphic
to a subgroup of \(G_i\).
\end{lemma}

\begin{proof}
Write \(G=\langle x_1,\ldots,x_m\mid r_1,\ldots,r_k\rangle\).  Choose
representatives of the finitely many generators in stages of the filtered
system and pass to one common later stage.  Each relator maps to the identity
in the colimit, so after passing to a further common later stage, all relators
are already trivial.  The chosen representatives therefore define a
homomorphism \(s:G\to G_i\), and \(\lambda_i\circ s\) fixes the generators of
the presented group \(G\).  Thus \(\lambda_i\circ s=\operatorname{id}_G\), so
\(s\) is injective.
\end{proof}

\begin{theorem}[Negative answer]\label{thm:negative}
The inductive-limit question in \cite{TianYu2026} has a negative answer:
\begin{enumerate}
  \item there is a countable finitely presented group that is not a directed
  union of groups with finite complexity rates;
  \item the same group is not a filtered colimit of groups with finite
  complexity rates, even when the bonding homomorphisms are arbitrary and need
  not be injective.
\end{enumerate}
\end{theorem}

\begin{proof}
Let \(G\) be the finitely presented group from Corollary~\ref{cor:fpbad}.

If \(G\) were a directed union of finite-complexity-rate subgroups, its finite
generating set would lie in one stage.  That stage would contain all of \(G\),
so \(G\) itself would have finite complexity rate, contradicting
\(\CR(G)=\infty\).

Now suppose \(G\cong\directlim_i G_i\), with all \(G_i\) of finite complexity
rate.  By Lemma~\ref{lem:retract}, \(G\) embeds in one \(G_i\).  Lemma
\ref{lem:subgroup} then implies finite complexity rate for \(G\), again a
contradiction.
\end{proof}

\section*{Verification Notes}

The proof is noncomputational.  The quantitative point to check is the
dependence on \(q\) in Lemma~\ref{lem:poincare}; the needed bound is linear in
\(q\), so with \(q=2p\) it gives the lower bound \(p\eps_p\ge c\).  The group
theoretic steps are formal once finite presentability and coarse expander
containment are available.

\section*{Novelty and Literature Status}

The exact notion of finite complexity rate was introduced in arXiv:2605.12930v1,
submitted May 13, 2026.  The previous run packet searched on June 16, 2026 for
the arXiv id, the paper title together with ``inductive limit'', and the phrase
``finite complexity rates inductive limit countable discrete group''; no later
answer to the exact question was found.  The supplied report records another
bounded search by title, arXiv id, the phrase ``finite complexity rate'', and the
inductive-limit question as of June 21, 2026, again without locating a later
explicit answer.  During this upgrade, the run indexes and web searches for the
arXiv id, finite complexity rate, and the supplied report title found no
duplicate packet or published answer.

This packet therefore records the result as a candidate counterexample produced
by synthesis of known ingredients: Matousek/Naor--Silberman expander Poincare
inequalities and Osajda's groups containing expanders.  It makes no priority
claim beyond that synthesis.

\section*{Human Review Recommendation}

Review as a candidate full negative answer to the source question.  The main
checks are:
\begin{enumerate}
  \item the source Definition 1.4 really implies \(p\eps_p\to 0\) for finite
  complexity rate;
  \item the cited Poincare inequality has the required linear \(q\)-dependence
  for \(\ell^q\);
  \item Osajda's Corollary 3.5 supplies a finitely presented group coarsely
  containing a bounded-degree expander sequence in the sense used above.
\end{enumerate}

\begin{thebibliography}{9}

\bibitem{Matousek1997}
Jiri Matousek,
\emph{On embedding expanders into \(\ell_p\) spaces},
Israel Journal of Mathematics \textbf{102} (1997), 189--197.

\bibitem{NaorSilberman2011}
Assaf Naor and Lior Silberman,
\emph{Poincare inequalities, embeddings, and wild groups},
Compositio Mathematica \textbf{147} (2011), no. 5, 1546--1572,
arXiv:1005.4084.

\bibitem{Osajda2020}
Damian Osajda,
\emph{Small cancellation labellings of some infinite graphs and applications},
Acta Mathematica \textbf{225} (2020), no. 1, 159--191,
arXiv:1406.5015.

\bibitem{TianYu2026}
Geng Tian and Guoliang Yu,
\emph{Embedding complexity into Banach spaces and the strong Novikov
conjecture},
arXiv:2605.12930v1, 2026.

\end{thebibliography}

\end{document}
